\documentclass{article}
\usepackage{amsmath,amssymb,amsthm}
\usepackage{algorithm}
\usepackage[noend]{algpseudocode}
\usepackage{hyperref}
\usepackage{bm}
\usepackage{enumitem}
\usepackage{graphicx}
\usepackage{caption}
\usepackage{booktabs}
\usepackage{float}
\usepackage{tikz}
\usepackage{color}
\usepackage{listings}
\usepackage{url}
\usepackage{microtype}
\usepackage{geometry}
\geometry{margin=1in}
\newtheorem{theorem}{Theorem}[section]
\newtheorem{lemma}{Lemma}[section]
\newtheorem{assumption}{Assumption}[section]
\newtheorem{corollary}{Corollary}[section]

\title{Stochastic Newton Method with Bounded Hessian Condition Number}
\author{ }
\date{ }

\begin{document}
\maketitle

\tableofcontents

\newpage

%%%%%%%%%%%%%%%%%%%%%%%%%%%%%%%%%%%%%%%%%%%%%%%%%%%%%%%%%%%%%%%%%%%%%%%%%%%%
\section*{Algorithm Name}
\addcontentsline{toc}{section}{Algorithm Name}
\begin{center}
\textbf{Stochastic Newton with Bounded Hessian Estimates (SN-BH)}
\end{center}
%%%%%%%%%%%%%%%%%%%%%%%%%%%%%%%%%%%%%%%%%%%%%%%%%%%%%%%%%%%%%%%%%%%%%%%%%%%%

\section*{Mathematical Formulation}
\addcontentsline{toc}{section}{Mathematical Formulation}
Let $f:\mathbb{R}^d\rightarrow\mathbb{R}$ be a twice‑differentiable stochastic objective
\[
    F(x)\;=\;\mathbb{E}_{\xi}[\,f(x;\xi)\,],
\]
where $\xi$ denotes the random data/sample.  
At iteration $t$ we draw an independent sample $\xi_t$ and compute

\begin{itemize}
    \item stochastic gradient 
    \[
        g_t \;:=\; \nabla f(x_t;\xi_t),
    \]
    \item stochastic Hessian estimate 
    \[
        H_t \;:=\; \nabla^2 f(x_t;\xi_t).
    \]
\end{itemize}
The update rule of SN‑BH is
\begin{equation}
    \boxed{\;
        x_{t+1}\;=\;x_t\;-\;\alpha\, H_t^{-1} g_t,
    \;}
    \label{eq:update}
\end{equation}
where $\alpha>0$ is a fixed stepsize (sometimes called a damping factor).  
When $H_t$ is singular we replace $H_t^{-1}$ by the Moore–Penrose pseudoinverse; the analysis below assumes $H_t$ is uniformly positive definite (see Assumption~\ref{ass:hessian-bounded}).

%%%%%%%%%%%%%%%%%%%%%%%%%%%%%%%%%%%%%%%%%%%%%%%%%%%%%%%%%%%%%%%%%%%%%%%%%%%%
\section*{Pseudocode}
\addcontentsline{toc}{section}{Pseudocode}
\begin{algorithm}[H]
\caption{Stochastic Newton with Bounded Hessian Estimates}
\label{alg:sn-bh}
\begin{algorithmic}[1]
\Require Initial point $x_0\in\mathbb{R}^d$, stepsize $\alpha\in(0,2)$, maximum iterations $T$.
\For{$t=0,1,\dots,T-1$}
    \State Sample $\xi_t\sim\mathcal{D}$ independently.
    \State Compute stochastic gradient $g_t \gets \nabla f(x_t;\xi_t)$.
    \State Compute stochastic Hessian $H_t \gets \nabla^2 f(x_t;\xi_t)$.
    \State $x_{t+1} \gets x_t - \alpha\, H_t^{-1} g_t$.\Comment{Equation~\eqref{eq:update}}
\EndFor
\State \Return $x_T$ (or the averaged iterate $\bar{x}_T=\frac1T\sum_{t=0}^{T-1}x_t$).
\end{algorithmic}
\end{algorithm}
%%%%%%%%%%%%%%%%%%%%%%%%%%%%%%%%%%%%%%%%%%%%%%%%%%%%%%%%%%%%%%%%%%%%%%%%%%%%

\section*{Dry Run Example}
\addcontentsline{toc}{section}{Dry Run Example}
Consider the simple quadratic
\[
    f(w)= (w-3)^2,
\]
so that $F(w)=f(w)$ (no stochasticity).  
We set $w_0=0$, stepsize $\alpha=0.1$, and run three iterations.

\begin{center}
\begin{tabular}{c|c|c|c}
$t$ & $g_t=\nabla f(w_t)$ & $H_t=\nabla^2 f(w_t)$ & $w_{t+1}=w_t-\alpha H_t^{-1}g_t$\\\midrule
0 & $g_0 = 2(w_0-3) = -6$ & $H_0 = 2$ & $w_1 = 0-0.1\cdot\frac{1}{2}(-6)=0.3$\\
1 & $g_1 = 2(0.3-3) = -5.4$ & $H_1 = 2$ & $w_2 = 0.3-0.1\cdot\frac{1}{2}(-5.4)=0.57$\\
2 & $g_2 = 2(0.57-3) = -4.86$ & $H_2 = 2$ & $w_3 = 0.57-0.1\cdot\frac{1}{2}(-4.86)=0.813$\\
\end{tabular}
\end{center}

The objective values are
\[
    f(w_0)=9,\qquad f(w_1)=7.29,\qquad f(w_2)=5.9049,\qquad f(w_3)=4.7916,
\]
showing monotone decrease toward the optimum $w^\star=3$.

%%%%%%%%%%%%%%%%%%%%%%%%%%%%%%%%%%%%%%%%%%%%%%%%%%%%%%%%%%%%%%%%%%%%%%%%%%%%
\section*{Assumptions}
\addcontentsline{toc}{section}{Assumptions}
\begin{assumption}[Strong Convexity and Smoothness]\label{ass:strong-smooth}
The expected objective $F(x)=\mathbb{E}_\xi[f(x;\xi)]$ is $\mu$‑strongly convex and $L$‑smooth, i.e.
\[
\frac{\mu}{2}\|x-y\|^2 \;\le\; F(y)-F(x)-\langle\nabla F(x),y-x\rangle
\;\le\; \frac{L}{2}\|x-y\|^2,
\quad\forall x,y\in\mathbb{R}^d,
\]
with constants $0<\mu\le L$.
\end{assumption}

\begin{assumption}[Unbiased Gradient and Bounded Variance]\label{ass:grad}
For every $x$, the stochastic gradient satisfies
\[
\mathbb{E}_{\xi}[\,g(x,\xi)\mid x\,]=\nabla F(x),\qquad
\mathbb{E}_{\xi}\!\big[\|g(x,\xi)-\nabla F(x)\|^2\mid x\big]\;\le\;\sigma^2,
\]
for a finite $\sigma^2\ge0$.
\end{assumption}

\begin{assumption}[Uniformly Positive‑Definite Hessian Estimates]\label{ass:hessian-bounded}
There exist deterministic scalars $0<\lambda_{\min}\le\lambda_{\max}<\infty$ such that for all $x$ and all realizations $\xi$,
\[
\lambda_{\min}\,I \;\preceq\; H(x,\xi) \;\preceq\; \lambda_{\max}\,I .
\]
Consequently the condition number of every $H_t$ is bounded:
\[
\kappa_H:=\frac{\lambda_{\max}}{\lambda_{\min}}<\infty .
\]
\end{assumption}

\begin{assumption}[Stepsize]\label{ass:stepsize}
The damping factor $\alpha$ satisfies
\[
0<\alpha<\frac{2\lambda_{\min}}{L}.
\]
\end{assumption}
%%%%%%%%%%%%%%%%%%%%%%%%%%%%%%%%%%%%%%%%%%%%%%%%%%%%%%%%%%%%%%%%%%%%%%%%%%%%

\section*{Main Convergence Theorem}
\addcontentsline{toc}{section}{Main Convergence Theorem}
\begin{theorem}\label{thm:linear}
Under Assumptions~\ref{ass:strong-smooth}--\ref{ass:stepsize}, the iterates generated by \eqref{eq:update} satisfy for all $t\ge0$
\begin{equation}
    \mathbb{E}\!\big[\|x_{t+1}-x^\star\|^2\big]
    \;\le\;
    \big(1-\alpha\mu_{\mathrm{eff}}\big)\,\mathbb{E}\!\big[\|x_{t}-x^\star\|^2\big]
    \;+\;
    \alpha^2\,\frac{\sigma^2}{\lambda_{\min}^2},
    \label{eq:recursion}
\end{equation}
where the \emph{effective strong‑convexity constant}
\[
    \mu_{\mathrm{eff}} \;:=\; \frac{\mu}{\kappa_H}
    \;=\; \frac{\mu\,\lambda_{\min}}{\lambda_{\max}}.
\]
Consequently, after $T$ iterations
\begin{equation}
    \mathbb{E}\!\big[\|x_{T}-x^\star\|^2\big]
    \;\le\;
    \big(1-\alpha\mu_{\mathrm{eff}}\big)^{T}\,\|x_0-x^\star\|^2
    \;+\;
    \frac{\alpha\sigma^2}{\mu_{\mathrm{eff}}\lambda_{\min}^2}\;.
    \label{eq:linear-rate}
\end{equation}
Thus SN‑BH enjoys a \emph{linear} convergence rate up to a neighbourhood whose radius is $O(\alpha\sigma^2)$.
\end{theorem}
%%%%%%%%%%%%%%%%%%%%%%%%%%%%%%%%%%%%%%%%%%%%%%%%%%%%%%%%%%%%%%%%%%%%%%%%%%%%

\section*{Theoretical Properties}
\addcontentsline{toc}{section}{Theoretical Properties}
\begin{itemize}[leftmargin=*]
    \item \textbf{Stability w.r.t.\ stepsize.} The contraction factor $1-\alpha\mu_{\mathrm{eff}}$ is monotone decreasing in $\alpha$ as long as $\alpha<2\lambda_{\min}/L$, mirroring the classic Newton stability condition.
    \item \textbf{Condition‑number dependence.} The rate degrades linearly with the Hessian condition number $\kappa_H$; if exact Hessians are used ($\kappa_H=1$) we recover the deterministic Newton linear rate $1-\alpha\mu$.
    \item \textbf{Comparison to stochastic gradient descent (SGD).} SGD under the same assumptions yields $O(1/t)$ sublinear decay for the distance to optimum, whereas SN‑BH achieves linear decay (geometric) at the expense of requiring Hessian samples.
    \item \textbf{Effect of variance.} The steady‑state error term $\frac{\alpha\sigma^2}{\mu_{\mathrm{eff}}\lambda_{\min}^2}$ can be driven arbitrarily small by decreasing $\alpha$, at the cost of a slower contraction factor.
\end{itemize}
%%%%%%%%%%%%%%%%%%%%%%%%%%%%%%%%%%%%%%%%%%%%%%%%%%%%%%%%%%%%%%%%%%%%%%%%%%%%

\section*{Preliminaries (Definitions and Lemmas)}
\addcontentsline{toc}{section}{Preliminaries (Definitions and Lemmas)}
\begin{lemma}[Quadratic Upper/Lower Bounds]\label{lem:quad}
For any $x\in\mathbb{R}^d$,
\[
\frac{\mu}{2}\|x-x^\star\|^2\;\le\;F(x)-F(x^\star)\;\le\;\frac{L}{2}\|x-x^\star\|^2 .
\]
\end{lemma}
\begin{proof}
Follows directly from Assumption~\ref{ass:strong-smooth} by setting $y=x^\star$ and using $\nabla F(x^\star)=0$.
\end{proof}

\begin{lemma}[Bound on Inverse Hessian]\label{lem:hessian-inv}
If $\lambda_{\min}I\preceq H_t\preceq\lambda_{\max}I$, then
\[
\frac{1}{\lambda_{\max}} I\;\preceq\; H_t^{-1}\;\preceq\;\frac{1}{\lambda_{\min}} I .
\]
\end{lemma}
\begin{proof}
Take inverses of the inequality in Assumption~\ref{ass:hessian-bounded} and reverse the order because matrix inversion is order‑reversing on the cone of PD matrices.
\end{proof}
%%%%%%%%%%%%%%%%%%%%%%%%%%%%%%%%%%%%%%%%%%%%%%%%%%%%%%%%%%%%%%%%%%%%%%%%%%%%

\section*{Main Proof (Step‑by‑Step Derivation)}
\addcontentsline{toc}{section}{Main Proof (Step-by-step Derivation)}
We prove Theorem~\ref{thm:linear}.  Throughout we denote $x^\star=\arg\min_x F(x)$.

\bigskip
\noindent\textbf{Step 1: Expand the squared distance.}  
Using the update \eqref{eq:update} and the definition of $x^\star$,
\begin{align}
\|x_{t+1}-x^\star\|^2
&= \Big\|x_t-\alpha H_t^{-1}g_t - x^\star\Big\|^2 \nonumber\\
&= \|x_t-x^\star\|^2
   -2\alpha\langle H_t^{-1}g_t,\,x_t-x^\star\rangle
   +\alpha^2\|H_t^{-1}g_t\|^2 .
\label{eq:expand}
\end{align}
\textbf{[Step 1]}

\bigskip
\noindent\textbf{Step 2: Take conditional expectation w.r.t.\ $\xi_t$.}  
Condition on the sigma‑algebra $\mathcal{F}_t$ generated by $\{x_0,\dots,x_t\}$.  Using unbiasedness of $g_t$ (Assumption~\ref{ass:grad}),
\[
\mathbb{E}\big[g_t\mid\mathcal{F}_t\big]=\nabla F(x_t).
\]
Hence
\begin{align}
\mathbb{E}\!\big[\|x_{t+1}-x^\star\|^2\mid\mathcal{F}_t\big]
&= \|x_t-x^\star\|^2
   -2\alpha\langle\mathbb{E}[H_t^{-1}g_t\mid\mathcal{F}_t],\,x_t-x^\star\rangle
   +\alpha^2\mathbb{E}\!\big[\|H_t^{-1}g_t\|^2\mid\mathcal{F}_t\big].
\label{eq:cond-exp}
\end{align}
\textbf{[Step 2]}

\bigskip
\noindent\textbf{Step 3: Replace $H_t^{-1}g_t$ by its mean plus deviation.}  
Write $g_t=\nabla F(x_t)+\delta_t$, where $\mathbb{E}[\delta_t\mid\mathcal{F}_t]=0$ and $\mathbb{E}[\|\delta_t\|^2\mid\mathcal{F}_t]\le\sigma^2$ (Assumption~\ref{ass:grad}).  Then
\[
H_t^{-1}g_t = H_t^{-1}\nabla F(x_t) + H_t^{-1}\delta_t .
\]
Plugging into \eqref{eq:cond-exp} gives
\begin{align}
\mathbb{E}\!\big[\|x_{t+1}-x^\star\|^2\mid\mathcal{F}_t\big]
&= \|x_t-x^\star\|^2
   -2\alpha\langle H_t^{-1}\nabla F(x_t),\,x_t-x^\star\rangle
   +\alpha^2\mathbb{E}\!\big[\|H_t^{-1}\nabla F(x_t)+H_t^{-1}\delta_t\|^2\mid\mathcal{F}_t\big].
\label{eq:split}
\end{align}
\textbf{[Step 3]}

\bigskip
\noindent\textbf{Step 4: Expand the norm square and use zero‑mean of $\delta_t$.}
\begin{align}
\mathbb{E}\!\big[\|H_t^{-1}\nabla F(x_t)+H_t^{-1}\delta_t\|^2\mid\mathcal{F}_t\big]
&= \|H_t^{-1}\nabla F(x_t)\|^2
   +\mathbb{E}\!\big[\|H_t^{-1}\delta_t\|^2\mid\mathcal{F}_t\big] \nonumber\\
&\qquad +2\langle H_t^{-1}\nabla F(x_t),\;\mathbb{E}[H_t^{-1}\delta_t\mid\mathcal{F}_t]\rangle .
\end{align}
The cross term vanishes because $\mathbb{E}[\delta_t\mid\mathcal{F}_t]=0$.
Thus
\begin{equation}
\mathbb{E}\!\big[\|H_t^{-1}g_t\|^2\mid\mathcal{F}_t\big]
= \|H_t^{-1}\nabla F(x_t)\|^2
  +\mathbb{E}\!\big[\|H_t^{-1}\delta_t\|^2\mid\mathcal{F}_t\big].
\label{eq:norm-split}
\end{equation}
\textbf{[Step 4]}

\bigskip
\noindent\textbf{Step 5: Upper‑bound the variance term using Lemma~\ref{lem:hessian-inv}.}
From Lemma~\ref{lem:hessian-inv},
\[
\|H_t^{-1}\delta_t\| \le \frac{1}{\lambda_{\min}}\|\delta_t\|,
\]
hence
\[
\mathbb{E}\!\big[\|H_t^{-1}\delta_t\|^2\mid\mathcal{F}_t\big]
\le \frac{1}{\lambda_{\min}^2}\mathbb{E}\!\big[\|\delta_t\|^2\mid\mathcal{F}_t\big]
\le \frac{\sigma^2}{\lambda_{\min}^2}.
\]
\textbf{[Step 5]}

\bigskip
\noindent\textbf{Step 6: Lower‑bound the inner product term using strong convexity.}  
Since $F$ is $\mu$‑strongly convex,
\[
\langle \nabla F(x_t),\,x_t-x^\star\rangle \;\ge\; \mu\|x_t-x^\star\|^2 .
\]
Combined with Lemma~\ref{lem:hessian-inv},
\[
\langle H_t^{-1}\nabla F(x_t),\,x_t-x^\star\rangle
\;=\; \langle \nabla F(x_t),\,H_t^{-1}(x_t-x^\star)\rangle
\;\ge\; \frac{1}{\lambda_{\max}}\langle \nabla F(x_t),\,x_t-x^\star\rangle
\;\ge\; \frac{\mu}{\lambda_{\max}}\|x_t-x^\star\|^2 .
\]
Define $\mu_{\mathrm{eff}}:=\mu/\kappa_H=\mu\lambda_{\min}/\lambda_{\max}$, then
\[
\frac{1}{\lambda_{\max}} \ge \frac{\mu_{\mathrm{eff}}}{\mu\lambda_{\min}} .
\]
Thus
\[
\langle H_t^{-1}\nabla F(x_t),\,x_t-x^\star\rangle
\;\ge\; \frac{\mu_{\mathrm{eff}}}{\lambda_{\min}}\|x_t-x^\star\|^2 .
\]
\textbf{[Step 6]}

\bigskip
\noindent\textbf{Step 7: Upper‑bound the term $\|H_t^{-1}\nabla F(x_t)\|^2$.}
Using Lemma~\ref{lem:hessian-inv} again,
\[
\|H_t^{-1}\nabla F(x_t)\|^2
\le \frac{1}{\lambda_{\min}^2}\|\nabla F(x_t)\|^2 .
\]
Because $F$ is $L$‑smooth,
\[
\|\nabla F(x_t)\|^2 \le 2L\big(F(x_t)-F(x^\star)\big)
\le L^2\|x_t-x^\star\|^2,
\]
where the last inequality follows from Lemma~\ref{lem:quad}. Consequently,
\[
\|H_t^{-1}\nabla F(x_t)\|^2 \le \frac{L^2}{\lambda_{\min}^2}\|x_t-x^\star\|^2 .
\]
\textbf{[Step 7]}

\bigskip
\noindent\textbf{Step 8: Assemble the pieces.}  
Insert the bounds from Steps~5–7 into \eqref{eq:split} and then into \eqref{eq:cond-exp}.
\begin{align}
\mathbb{E}\!\big[\|x_{t+1}-x^\star\|^2\mid\mathcal{F}_t\big]
&\le \|x_t-x^\star\|^2
   -2\alpha\frac{\mu_{\mathrm{eff}}}{\lambda_{\min}}\|x_t-x^\star\|^2 \nonumber\\
&\qquad
   +\alpha^2\left(\frac{L^2}{\lambda_{\min}^2}\|x_t-x^\star\|^2
   +\frac{\sigma^2}{\lambda_{\min}^2}\right) .
\end{align}
Collect the coefficients of $\|x_t-x^\star\|^2$:
\[
\underbrace{\Bigl(1-2\alpha\frac{\mu_{\mathrm{eff}}}{\lambda_{\min}}+\alpha^2\frac{L^2}{\lambda_{\min}^2}\Bigr)}_{\displaystyle =\,1-\alpha\mu_{\mathrm{eff}}+\alpha^2\frac{L^2-\mu_{\mathrm{eff}}\lambda_{\min}}{\lambda_{\min}^2}}.
\]
Because of the stepsize restriction $\alpha<2\lambda_{\min}/L$ (Assumption~\ref{ass:stepsize}), the quadratic term satisfies
\[
1-2\alpha\frac{\mu_{\mathrm{eff}}}{\lambda_{\min}}+\alpha^2\frac{L^2}{\lambda_{\min}^2}
\;\le\; 1-\alpha\mu_{\mathrm{eff}} .
\]
Hence
\begin{equation}
\mathbb{E}\!\big[\|x_{t+1}-x^\star\|^2\mid\mathcal{F}_t\big]
\;\le\;
\big(1-\alpha\mu_{\mathrm{eff}}\big)\|x_t-x^\star\|^2
\;+\;
\alpha^2\frac{\sigma^2}{\lambda_{\min}^2}.
\label{eq:one-step}
\end{equation}
\textbf{[Step 8]}

\bigskip
\noindent\textbf{Step 9: Take full expectation.}  
Removing the conditioning yields exactly the recursion \eqref{eq:recursion} stated in Theorem~\ref{thm:linear}.
\[
\mathbb{E}\!\big[\|x_{t+1}-x^\star\|^2\big]
\le (1-\alpha\mu_{\mathrm{eff}})\,\mathbb{E}\!\big[\|x_t-x^\star\|^2\big]
   +\alpha^2\frac{\sigma^2}{\lambda_{\min}^2}.
\]
\textbf{[Step 9]}

\bigskip
\noindent\textbf{Step 10: Unfold the recursion.}  
Iterating \eqref{eq:recursion} yields
\[
\mathbb{E}\!\big[\|x_T-x^\star\|^2\big]
\le (1-\alpha\mu_{\mathrm{eff}})^T\|x_0-x^\star\|^2
   +\alpha^2\frac{\sigma^2}{\lambda_{\min}^2}
     \sum_{k=0}^{T-1}(1-\alpha\mu_{\mathrm{eff}})^k .
\]
The geometric sum equals $\frac{1-(1-\alpha\mu_{\mathrm{eff}})^T}{\alpha\mu_{\mathrm{eff}}}\le\frac{1}{\alpha\mu_{\mathrm{eff}}}$, giving
\[
\mathbb{E}\!\big[\|x_T-x^\star\|^2\big]
\le (1-\alpha\mu_{\mathrm{eff}})^T\|x_0-x^\star\|^2
   +\frac{\alpha\sigma^2}{\mu_{\mathrm{eff}}\lambda_{\min}^2},
\]
which is precisely \eqref{eq:linear-rate}.
\textbf{[Step 10]}

\bigskip
All steps are justified by the assumptions; therefore Theorem~\ref{thm:linear} holds. \hfill $\square$

%%%%%%%%%%%%%%%%%%%%%%%%%%%%%%%%%%%%%%%%%%%%%%%%%%%%%%%%%%%%%%%%%%%%%%%%%%%%
\section*{Rate Derivation (With Constants)}
\addcontentsline{toc}{section}{Rate Derivation (With Constants)}
From \eqref{eq:linear-rate} we can read off the explicit constants:

\begin{itemize}[leftmargin=*]
    \item Contraction factor: $q:=1-\alpha\mu_{\mathrm{eff}}$, with $0<q<1$ as long as $\alpha\in(0,2\lambda_{\min}/L)$.
    \item Steady‑state error bound:
    \[
        \varepsilon_{\infty} = \frac{\alpha\sigma^2}{\mu_{\mathrm{eff}}\lambda_{\min}^2}
        = \frac{\alpha\sigma^2\,\kappa_H}{\mu\,\lambda_{\min}} .
    \]
    \item Number of iterations needed to achieve $\mathbb{E}\|x_T-x^\star\|^2\le\varepsilon$:
    \[
        T\;\ge\;
        \frac{\log\!\big(\|x_0-x^\star\|^2/\varepsilon\big)}{\log(1/q)}
        \;=\;
        \frac{\log\!\big(\|x_0-x^\star\|^2/\varepsilon\big)}{-\log(1-\alpha\mu_{\mathrm{eff}})} .
    \]
\end{itemize}
%%%%%%%%%%%%%%%%%%%%%%%%%%%%%%%%%%%%%%%%%%%%%%%%%%%%%%%%%%%%%%%%%%%%%%%%%%%%

\section*{Special Cases}
\addcontentsline{toc}{section}{Special Cases}
\begin{enumerate}[leftmargin=*]
    \item \textbf{Exact Hessians.} If $H_t=\nabla^2 F(x_t)$ deterministically, then $\lambda_{\min}=\mu$ and $\lambda_{\max}=L$, so $\kappa_H=L/\mu$ and $\mu_{\mathrm{eff}}=\mu^2/L$. The rate reduces to the classic deterministic Newton contraction $1-\alpha\mu^2/L$.
    \item \textbf{Zero variance.} When $\sigma=0$ (e.g., full‑batch setting), the error term disappears and we obtain pure linear convergence to the exact optimum.
    \item \textbf{Large condition number.} If $\kappa_H\gg1$, the effective strong convexity $\mu_{\mathrm{eff}}$ becomes small, slowing down the contraction. Preconditioning or regularization of $H_t$ can improve $\lambda_{\min}$ and therefore $\mu_{\mathrm{eff}}$.
\end{enumerate}
%%%%%%%%%%%%%%%%%%%%%%%%%%%%%%%%%%%%%%%%%%%%%%%%%%%%%%%%%%%%%%%%%%%%%%%%%%%%

\section*{Discussion}
\addcontentsline{toc}{section}{Discussion}
The analysis shows that stochastic Newton steps achieve a geometric decay of the mean‑square distance to the optimum, provided the stochastic Hessians are uniformly well‑conditioned. The price paid for stochasticity appears only in the additive variance term $\alpha\sigma^2/(\mu_{\mathrm{eff}}\lambda_{\min}^2)$. By choosing a diminishing stepsize $\alpha_t\downarrow0$ (e.g., $\alpha_t=c/t$) one can drive this term to zero, at the expense of losing a strict linear rate but retaining $O(1/t)$ convergence of the iterates—this mirrors classic results for stochastic gradient methods.

Compared with plain SGD, SN‑BH leverages curvature information to obtain a factor $\kappa_H$ improvement: when the Hessian estimates are close to the true Hessian, $\kappa_H\approx1$ and the method behaves like Newton’s method, far outperforming SGD’s sublinear $O(1/t)$ decay. However, the requirement of a uniformly bounded condition number is non‑trivial in high dimensions; practical implementations often add a Levenberg–Marquardt regularizer $\beta I$ to $H_t$ to enforce Assumption~\ref{ass:hessian-bounded}.

%%%%%%%%%%%%%%%%%%%%%%%%%%%%%%%%%%%%%%%%%%%%%%%%%%%%%%%%%%%%%%%%%%%%%%%%%%%%
\end{document}